\input{../Headers/Preamble.tex}
\usepackage[utf8]{inputenc}
\usepackage[letterpaper, margin=1in, bottom=2.5in, top=2.2in]{geometry}
\usepackage{fancyhdr}
\setlength{\headheight}{15pt}
\setlength{\footskip}{72pt}

\usepackage{xcolor}
\usepackage{amsmath}
\usepackage{anyfontsize}
\usepackage{graphicx}
\usepackage{tikz}
\usepackage{listings}
\usepackage{sourcesanspro}
\renewcommand{\familydefault}{\sfdefault}
\usepackage{fontawesome5}
\usepackage{titlesec}
\usepackage{setspace}
\usepackage{hyperref}
\usepackage{mdframed}
\usepackage{qrcode} % Añadir paquete para generar códigos QR
\usetikzlibrary{calc,shapes,positioning}
\usepackage{eso-pic}

\AtBeginDocument{\color{primaryColor}}

% Colors updated with JavaScript brand colors
\definecolor{bgColor}{RGB}{13, 17, 23}  % Dark background similar to JS docs
\definecolor{primaryColor}{RGB}{255, 255, 255}  % White text for better contrast
\definecolor{accentColor}{RGB}{255, 220, 60}  % Brighter JavaScript Yellow for better contrast
\definecolor{jsSecondary}{RGB}{180, 180, 180}  % Light gray as secondary color for better contrast
\definecolor{secondaryColor}{RGB}{230, 235, 240}  % Almost white gray for better contrast
\definecolor{terminalBg}{RGB}{22, 22, 22}  % Terminal background
\definecolor{terminalFrame}{RGB}{40, 40, 40}  % Terminal frame
\definecolor{lineNumberColor}{RGB}{100, 100, 100}  % Line number color
\definecolor{dividerColor}{RGB}{120, 130, 150}  % Brighter divider color for better visibility

% Colores para código 
\definecolor{codeTextColor}{RGB}{255, 255, 255}  % Blanco puro para el texto básico
\definecolor{codeKeywordColor}{RGB}{97, 175, 239}  % Azul claro para palabras clave
\definecolor{codeCommentColor}{RGB}{121, 192, 120}  % Verde para comentarios
\definecolor{codeStringColor}{RGB}{255, 220, 60}  % Amarillo JS más brillante para strings
\definecolor{codeMethodColor}{RGB}{255, 160, 122}  % Naranja claro para métodos
\definecolor{codeFunctionColor}{RGB}{230, 180, 80}  % Tono amarillo-dorado para funciones
\definecolor{codeNumberColor}{RGB}{180, 230, 180}  % Verde claro para números

\pagecolor{bgColor}

\hypersetup{
    colorlinks=true,
    linkcolor=accentColor,
    filecolor=accentColor,
    urlcolor=accentColor,
}

% Definir lenguaje JavaScript para listings
\lstdefinelanguage{JavaScript}{
  keywords={break, case, catch, continue, debugger, default, delete, do, else, false, finally, for, function, if, in, instanceof, new, null, return, switch, this, throw, true, try, typeof, var, void, while, with, let, const, class, export, import, yield, async, await},
  sensitive=true,
  comment=[l]{//},
  morecomment=[s]{/*}{*/},
  morestring=[b]',
  morestring=[b]"
}

% Configuración para listings con JavaScript como lenguaje predeterminado
\lstset{
  language=JavaScript,
  basicstyle=\ttfamily\bfseries\color{codeTextColor},
  backgroundcolor=\color{terminalBg},
  commentstyle=\color{codeCommentColor},
  keywordstyle=\color{codeKeywordColor},
  stringstyle=\color{codeStringColor},
  numberstyle=\color{lineNumberColor},
  breaklines=true,
  breakatwhitespace=true,
  tabsize=2,
  showstringspaces=false,
  frame=none,
  xleftmargin=15pt,
  xrightmargin=0pt,
  aboveskip=5pt,
  belowskip=5pt,
  numbers=left,
  numbersep=8pt,
  extendedchars=true,
  keepspaces=true,
  columns=flexible,
  lineskip=2pt,
  emph={[2]require,createReadStream,createWriteStream,pipe,on,once,createGzip,pipeline,Transform,Readable,Writable,Duplex},
  emphstyle={[2]\color{codeFunctionColor}}
}

\newenvironment{macterminal}{%
    \begin{mdframed}[
        linecolor=terminalFrame,
        backgroundcolor=terminalBg,
        roundcorner=5pt,
        skipabove=5pt,
        skipbelow=5pt,
        linewidth=1pt,
        innertopmargin=5pt,
        frametitle={%
            \tikz[baseline=(current bounding box.east), outer sep=0pt]{
                \fill[red!80!black] (0,0) circle (5pt);
                \fill[yellow!80!black] (0.7,0) circle (5pt);
                \fill[green!70!black] (1.4,0) circle (5pt);
            }
        },
        frametitlealignment=\raggedright,
        frametitleaboveskip=8pt,
        frametitlebelowskip=0pt,
    ]
}{%
    \end{mdframed}%
}

\newcommand{\verspace}{\vspace{5pt}}

\titleformat{\section}
  {\LARGE\bfseries\color{primaryColor}}
  {\thesection. }
  {0pt}
  {}
  []

\titleformat{\subsection}
  {\Large\bfseries\color{accentColor}}
  {\thesubsection. }
  {0pt}
  {}
  []

\titleformat{\subsubsection}
  {\large\bfseries\color{primaryColor}}
  {\thesubsubsection. }
  {0pt}
  {}
  []

\titlespacing*{\section}{0pt}{18pt}{8pt}
\titlespacing*{\subsection}{0pt}{12pt}{5pt}
\titlespacing*{\subsubsection}{0pt}{8pt}{3pt}

\pagestyle{fancy}
\fancyhf{}

\fancyhead[L]{
    \begin{tikzpicture}[remember picture, overlay]
        \fill[accentColor, rounded corners=3pt] (0,0) rectangle (2.2cm,0.7cm);
        \node[text=bgColor, font=\bfseries] at (1.1cm,0.35cm) {NODE.JS};
        \fill[secondaryColor] (2.35cm,0.1cm) rectangle (2.40cm,0.6cm);
        \node[text=primaryColor, font=\bfseries, anchor=west] at (2.35cm,0.35cm) {STREAMS};
    \end{tikzpicture}
}

\renewcommand{\headrulewidth}{0pt}
\renewcommand{\footrulewidth}{0.2pt}
\renewcommand{\footruleskip}{1cm}

\fancyfoot[C]{
    \vspace*{0.1cm}
    \noindent
    \begin{minipage}{\textwidth}
        \begin{flushleft}
            % Foto de perfil
            \raisebox{0.7cm}{
            \begin{tikzpicture}[baseline]
                \path[fill=bgColor] (0,0) circle (0.6cm);
                \clip (0,0) circle (0.6cm);
                \node at (0,0) {
                    \includegraphics[width=1.2cm,height=1.2cm]{../../Images/profile-image.jpeg}
                };
            \end{tikzpicture}
            }
            % Información de perfil
            \begin{minipage}[b]{0.5\textwidth}
                {\large\bfseries\color{primaryColor}Alejandro Sánchez Yalí}
                \par\vspace{1pt}
                {\small\color{secondaryColor}Software Developer | AI \& Blockchain Enthusiast}
                \par\vspace{1pt}
                {\small\color{accentColor}\faGlobe\hspace{5pt}\color{secondaryColor}www.asanchezyali.com}
            \end{minipage}
        \end{flushleft}
        \vspace{6pt}
    \end{minipage}
}

\fancypagestyle{plain}{
    \fancyhf{}
    \renewcommand{\headrulewidth}{0pt}
    \renewcommand{\footrulewidth}{0pt}
}

\renewcommand{\labelitemi}{\textcolor{accentColor}{$\bullet$}}
\renewcommand{\labelitemii}{\textcolor{secondaryColor}{$\circ$}}

\usepackage{relsize}
\AtBeginDocument{\relsize{1}}

\newcommand{\languagetag}[1]{
    \begin{tikzpicture}[baseline]
        \node[fill=accentColor, text=bgColor, rounded corners=5pt, inner sep=7pt] {
            {\normalsize\textbf{#1}}
        };
    \end{tikzpicture}
}

\newcommand{\BackgroundPic}{%
    \put(0,0){%
        \begin{tikzpicture}[remember picture, overlay]
            % Coloca la imagen de fondo que cubra toda la página
            \node[inner sep=0pt] at (current page.center) {
                \includegraphics[width=\paperwidth,height=\paperheight,keepaspectratio=false]{../../Images/Intense Focus at Gaming Event.jpeg}
            };
            
            % Agrega una capa semitransparente oscura sobre la imagen
            % Puedes ajustar el color y la opacidad según necesites
            \fill[black, opacity=0.7] (current page.south west) rectangle (current page.north east);
            
            % Si quieres un degradado oscuro en lugar de un color sólido, usa esto:
            % \shade[top color=black!80, bottom color=black!40, opacity=0.7] 
            %    (current page.south west) rectangle (current page.north east);
        \end{tikzpicture}%
    }%
}

% Comando para QR elegante (para la página de título)
\newcommand{\elegantqr}[2]{
    \qrcode[height=2.5cm]{#1}
    \\[0.1cm]
    {\hspace{0.2cm}\color{primaryColor}\small #2\par}
}

\newcommand{\titlepagecontents}{%
    \AddToShipoutPicture*{\BackgroundPic}
    \vspace*{3cm}
    % Asegurar que todos los elementos estén alineados a la izquierda
    \begin{flushleft}
    \languagetag{Node.js}\\[0.4cm]
    {\fontsize{46}{52}\bfseries\color{primaryColor}Parallelization \& \color{accentColor}Distributed \color{primaryColor}Computing\par}
    \vspace{0.3cm}
    {\fontsize{18}{52}\color{secondaryColor}Unleashing Multi-Core Power in Node.js\par} 
    \vspace{0.3cm}
    {\color{secondaryColor}\today\par}
    \vspace{2cm}
    % El QR se alineará a la izquierda aunque internamente esté centrado
    \elegantqr{https://github.com/asanchezyali/social-media-posts/tree/main/JavaScript/Parallelization}{Source Code}
    \end{flushleft}
}

\newcommand{\finalpagecontents}{%
    \AddToShipoutPicture*{\BackgroundPic}
    \vspace*{3cm}
    % Left-aligned elements with EXACT same structure as title page
    \begin{flushleft}
    \languagetag{Feedback}\\[0.4cm]
    {\fontsize{46}{52}\bfseries\color{primaryColor}Found this \color{accentColor}helpful?\par}
    \vspace{0.3cm}
    {\fontsize{18}{52}\color{secondaryColor}Save, comment and share\par}  
    \vspace{0.3cm}
    {\color{secondaryColor}\today\par}
    \end{flushleft}
}

\begin{document}
\setlength{\parindent}{0pt}
\color{primaryColor}

\begin{titlepage}
    \titlepagecontents
\end{titlepage}

\section{Breaking the Single-Thread Barrier in Node.js}

Node.js is renowned for its non-blocking, event-driven architecture built on a single-threaded execution model. Although Node.js excels at handling concurrent I/O operations thanks to its event model, CPU-intensive tasks can block the event loop, which degrades application performance.

Parallel and distributed computing techniques offer three key advantages:

\begin{itemize}
    \item \textbf{\textcolor{accentColor}{Enhanced Performance:}} Distribute workloads across multiple cores and machines.
    
    \item \textbf{\textcolor{accentColor}{Event Loop Protection:}} Prevent CPU-bound tasks from blocking the main thread.
    
    \item \textbf{\textcolor{accentColor}{Horizontal Scalability:}} Scale beyond the limitations of a single machine.
\end{itemize}

\subsection{Understanding Node.js Execution Model}

Node.js operates on a single-threaded event loop, which excels at handling concurrent I/O but struggles with CPU-intensive tasks:

\begin{macterminal}
\begin{lstlisting}
// CPU-intensive operation blocking the event loop
const http = require("http");

// This function performs an intensive operation that will block the event loop
function computeIntensive() {
  let result = 0;
  for (let i = 0; i < 1e9; i++) {
    result += i;
  }
  return result;
}

const server = http.createServer((req, res) => {
  if (req.url === "/compute" && req.method === "GET") {
    // This call completely blocks the server
    // while the calculation executes
    const result = computeIntensive();
    res.writeHead(200, { "Content-Type": "application/json" });
    res.end(JSON.stringify({ result }));
  } else {
    res.writeHead(404, { "Content-Type": "text/plain" });
    res.end("Not Found");
  }
});

const PORT = 3000;
server.listen(PORT, () => {
  console.log(`Server running at http://localhost:${PORT}`);
});
\end{lstlisting}
\end{macterminal}

This code demonstrates the fundamental problem: when a CPU-intensive operation runs in Node.js, it completely blocks the event loop. The loop in \textbf{\textcolor{accentColor}{computeIntensive()}} performs a simple but heavy task that prevents any other request from being served until it completes. In production applications, this would cause timeouts and poor user experience.

This causes significant limitations:
\begin{itemize}
    \item CPU-bound tasks block the entire application
    \item Only a single core is utilized regardless of the machine's capabilities
\end{itemize}

\section{Parallel Computing with Worker Threads}

Worker Threads provide true parallelism by enabling the creation of separate JavaScript execution threads. Unlike the cluster module, which is more suitable for I/O-bound tasks, Worker Threads are specifically designed for CPU-intensive tasks by allowing JavaScript code execution in separate threads.

\begin{macterminal}
\begin{lstlisting}
// main.js - Parent thread implementation
const { Worker } = require('worker_threads');
const os = require('os');

// Determine available CPU cores
const numCPUs = os.cpus().length;

// This function creates and manages a worker thread
function runWorker(workerData) {
  return new Promise((resolve, reject) => {
    const worker = new Worker('./worker.js', { workerData });
    worker.on('message', resolve);
    worker.on('error', reject);
    worker.on('exit', (code) => {
      if (code !== 0) reject(new Error(`Worker stopped with code ${code}`));
    });
  });
}

async function main() {
  try {
    console.time('parallel-execution');
    
    // Create a worker for each available CPU core
    // and divide the work into equal parts
    const workers = Array(numCPUs).fill().map((_, index) => {
      return runWorker({
        start: index * 250000000,  // Divide the total range among the cores
        end: (index + 1) * 250000000
      });
    });
    
    // Wait for all workers to complete their work
    const results = await Promise.all(workers);
    // Combine the partial results
    const finalResult = results.reduce((sum, current) => sum + current, 0);
    
    console.timeEnd('parallel-execution');
    console.log(`Final sum: ${finalResult}`);
  } catch (err) {
    console.error('Error:', err);
  }
}

main();

// Output
// parallel-execution: 352.207ms
// Final sum: 3124999997616684500
\end{lstlisting}
\end{macterminal}

This code implements the solution using Worker Threads. The main thread:
\begin{itemize}
    \item Automatically detects the number of available cores with \textbf{\textcolor{accentColor}{os.cpus().length}}
    \item Encapsulates each worker in a promise to manage its lifecycle
    \item Divides the work equally for each core
    \item Uses \textbf{\textcolor{accentColor}{Promise.all}} to wait for all workers to finish
    \item Combines the partial results to get the final result
\end{itemize}

The worker thread contains the CPU-intensive code:

\begin{macterminal}
\begin{lstlisting}
// worker.js - The code executed in parallel
const { parentPort, workerData } = require('worker_threads');

// This function performs the intensive calculation
// but only in the range assigned to this worker
function computeIntensiveTask(start, end) {
  let sum = 0;
  for (let i = start; i < end; i++) {
    sum += i;
  }
  return sum;
}

// Extract the range limits assigned to this worker
const { start, end } = workerData;
// Execute the calculation only in this range
const result = computeIntensiveTask(start, end);

// Send the result back to the main thread
parentPort.postMessage(result);
\end{lstlisting}
\end{macterminal}

The worker code is simpler. Each worker:
\begin{itemize}
    \item Receives its portion of work through \textbf{\textcolor{accentColor}{workerData}}
    \item Executes the CPU-intensive task only in its assigned range
    \item Communicates the result back to the main thread using \textbf{\textcolor{accentColor}{parentPort.postMessage()}}
\end{itemize}

The key is that multiple instances of this code will run simultaneously in different threads, taking advantage of all available cores without blocking the main event loop.

\section{Distributed Computing with Node.js}

Distributed computing expands computational capacity across multiple machines for virtually unlimited scaling. Unlike parallel computing that operates on a single machine, distributed computing allows workload distribution among multiple servers, offering greater scalability and fault tolerance.

\subsection{Message Queue Architecture}

A message queue forms the backbone of distributed systems:

\begin{macterminal}
\begin{lstlisting}
// Bull is a popular library for managing queues in Node.js
// There are other alternatives like RabbitMQ or Kafka
const Queue = require("bull");

const workQueue = new Queue("data-processing", {
  redis: { host: "localhost", port: 6379 },
});

async function distributeWork(dataset, chunkSize = 1000) {
  // Split the dataset into smaller chunks
  const chunks = [];
  for (let i = 0; i < dataset.length; i += chunkSize) {
    chunks.push(dataset.slice(i, i + chunkSize));
  }

  console.log(`Distributing ${chunks.length} tasks`);

  // Create a job in the queue for each data chunk
  const jobPromises = chunks.map((chunk, index) => {
    return workQueue.add(
      "process-chunk",
      {
        chunkId: index,
        data: chunk,
        timestamp: Date.now(),
      },
      {
        attempts: 3,  // Automatic retry if the job fails
      }
    );
  });

  await Promise.all(jobPromises);
  console.log("All tasks queued for processing");
}

// Simulate a large dataset
// This could be any data source, like a database or an API
const mockData = Array.from({ length: 10000 }, (_, i) => ({ id: i, value: Math.random() }));

distributeWork(mockData, 1000);
\end{lstlisting}
\end{macterminal}

This code implements the producer in a distributed queue system:
\begin{itemize}
    \item Uses \textbf{\textcolor{accentColor}{Bull}}, a Redis-based library for queue management
    \item Divides a large dataset into smaller chunks
    \item Adds each chunk as an independent job to the queue
    \item Configures automatic retries to ensure job completion
    \item Operates asynchronously to avoid blocking the service while queueing tasks
\end{itemize}

This approach allows multiple servers or processes to pick up and process jobs independently.

\subsection{Worker Process Implementation}

Worker processes run on separate machines. To achieve true distributed computing, this code should be executed on different servers with a shared Redis instance configured to be accessible from all machines in the network.

\begin{macterminal}
\begin{lstlisting}
// Consumer: Process tasks from the queue
const Queue = require('bull');
const cluster = require('cluster');
const os = require('os');

// Note: In a real distributed environment, this host would point
// to a Redis instance accessible from all machines
const workQueue = new Queue('data-processing', {
  redis: { host: 'redis-server', port: 6379 }
});

// Use cluster to utilize all CPU cores
if (cluster.isMaster) {
  const numCPUs = os.cpus().length;
  console.log(`Setting up ${numCPUs} workers`);
  
  // Create a worker process for each CPU
  for (let i = 0; i < numCPUs; i++) {
    cluster.fork();
  }
  
  // Implement fault tolerance by automatically restarting failed workers
  cluster.on('exit', (worker) => {
    console.log(`Worker ${worker.process.pid} died. Restarting...`);
    cluster.fork();
  });
} else {
  // Worker process logic
  console.log(`Worker ${process.pid} started`);
  
  // Process jobs (8 concurrent jobs per worker)
  workQueue.process('process-chunk', 8, async (job) => {
    const { chunkId, data } = job.data;
    await job.progress(10);  // Report initial progress
    
    console.log(`Processing chunk ${chunkId}`);
    const startTime = Date.now();
    
    // Process data (CPU-intensive work)
    const processedData = data.map(item => {
      let result = 0;
      for (let i = 0; i < 10000; i++) {
        result += Math.sqrt(item.value * i);
      }
      return { ...item, processed: result };
    });
    
    await job.progress(100);  // Report completion
    const duration = Date.now() - startTime;
    
    return {
      chunkId,
      processedCount: processedData.length,
      duration,
      workerId: process.pid
    };
  });
}
// Output
// Setting up 10 workers
// Worker 11419 started
// Worker 11421 started
// Worker 11426 started
// Worker 11423 started
// Worker 11420 started
// Worker 11422 started
// Worker 11418 started
// Worker 11425 started
// Worker 11427 started
// Worker 11424 started
\end{lstlisting}
\end{macterminal}

This code implements the consumer in a distributed queue system:
\begin{itemize}
    \item Uses the \textbf{\textcolor{accentColor}{cluster}} module to create multiple processes within each machine
    \item The main (master) process manages the creation and supervision of workers
    \item Implements fault tolerance by automatically restarting failed processes
    \item Worker processes consume jobs from the shared Redis queue
    \item Each worker handles up to 8 concurrent jobs, balancing parallelism and resources
    \item Reports progress during processing to allow real-time monitoring
\end{itemize}

This combination of cluster and distributed queues enables scaling processing both vertically (more cores) and horizontally (more machines).

\section{Advanced Patterns for Parallel Processing}

\subsection{Worker Pool Pattern}

For frequent parallel operations, a worker pool provides better efficiency:

\begin{macterminal}
\begin{lstlisting}
const { Worker } = require('worker_threads');
const os = require('os');

class WorkerPool {
  constructor(workerScript, numWorkers = os.cpus().length) {
    this.workerScript = workerScript;
    this.numWorkers = numWorkers;
    this.workers = [];
    this.freeWorkers = [];
    this.taskQueue = [];
    this.tasks = {};
    this.taskIdCounter = 0;
    
    this._initPool();
  }
  
  // Initialize the worker pool
  _initPool() {
    // Create the worker threads
    for (let i = 0; i < this.numWorkers; i++) {
      const worker = new Worker(this.workerScript);
      this.workers.push(worker);
      this.freeWorkers.push(i);
      
      // Set up message handling
      worker.on('message', (result) => {
        const { taskId, data } = result;
        
        // Resolve the promise for this task
        if (this.tasks[taskId]) {
          this.tasks[taskId].resolve(data);
          delete this.tasks[taskId];
        }
        
        // Return the worker to the available pool
        this.freeWorkers.push(i);
        this._processQueue();
      });
    }
  }
  
  // Public API to execute tasks
  executeTask(taskData) {
    return new Promise((resolve, reject) => {
      const taskId = this.taskIdCounter++;
      
      this.tasks[taskId] = { resolve, reject, taskData };
      this.taskQueue.push(taskId);
      this._processQueue();
    });
  }
  
  // Process task queue if workers are available
  _processQueue() {
    if (this.taskQueue.length > 0 && this.freeWorkers.length > 0) {
      const taskId = this.taskQueue.shift();
      const workerId = this.freeWorkers.shift();
      
      this.workers[workerId].postMessage({
        taskId,
        data: this.tasks[taskId].taskData
      });
    }
  }
}

// Example of using the WorkerPool
async function example() {
  // Create a pool with a specific worker script
  const pool = new WorkerPool('./calculation-worker.js');
  
  // Execute multiple tasks in parallel
  const results = await Promise.all([
    pool.executeTask({ type: 'sum', values: [1, 2, 3, 4, 5] }),
    pool.executeTask({ type: 'product', values: [2, 3, 4] }),
    pool.executeTask({ type: 'statistics', values: [10, 20, 30, 40, 50] })
  ]);
  
  console.log(results);
}
\end{lstlisting}
\end{macterminal}

This Worker Pool pattern optimizes the use of worker threads:
\begin{itemize}
    \item Creates a fixed set of workers at startup, avoiding the overhead of creating/destroying threads
    \item Implements a queuing system for incoming tasks
    \item Automatically manages task assignment to available workers
    \item Provides a promise-based interface for executing tasks in parallel
    \item Reuses workers after completing tasks for optimal performance
\end{itemize}

The WorkerPool is ideal for applications that require frequent parallel operations, such as servers that process images, analyze data, or perform calculations for multiple simultaneous users.

\section{Best Practices for Node.js Parallelization}

\begin{itemize}
    \item \textbf{\textcolor{accentColor}{Task Granularity:}} Divide work into optimally sized chunks that balance overhead and parallelism.
    Example: For image processing, divide into blocks of 100-1000 images instead of processing one by one.
    
    \item \textbf{\textcolor{accentColor}{Resource Management:}} Monitor memory and CPU usage to prevent overloading systems.
    Example: Implement concurrency limits based on available memory and use tools like \texttt{node-memwatch}.
    
    \item \textbf{\textcolor{accentColor}{Error Handling:}} Implement comprehensive error recovery strategies in distributed systems.
    Example: Combine exponential backoff retries, circuit breakers, and dead-letter queues to manage failures.
    
    \item \textbf{\textcolor{accentColor}{Data Serialization:}} Minimize data transfer between threads and processes for better performance.
    Example: Transfer only identifiers or references when possible instead of complete objects.
    
    \item \textbf{\textcolor{accentColor}{Backpressure:}} Implement flow control mechanisms to prevent overwhelming workers.
    Example: Limit the number of tasks that can be in queue or use rate limiting systems.
\end{itemize}

\section{When To Use Each Approach}

Below is a comparison of the different approaches to help choose the most appropriate technique according to the use case:

\begin{center}
\small  % o \footnotesize para hacerla aún más pequeña
\begin{tabular}{|p{3cm}|p{3.5cm}|p{3.5cm}|p{3.5cm}|}
\hline
\textcolor{accentColor}{\textbf{Approach}} & \textcolor{primaryColor}{\textbf{Strengths}} & \textcolor{primaryColor}{\textbf{Limitations}} & \textcolor{primaryColor}{\textbf{Use Cases}} \\
\hline
\textbf{\textcolor{accentColor}{Worker Threads}} & 
True parallelism, efficient memory sharing & 
Limited to a single machine & 
CPU-intensive tasks (calculations, data processing) \\
\hline
\textbf{\textcolor{accentColor}{Cluster}} & 
Easy to implement, good load balancing & 
Not optimal for CPU-intensive tasks & 
Web servers, APIs, I/O-bound applications \\
\hline
\textbf{\textcolor{accentColor}{Distributed Queues}} & 
Horizontal scalability, fault tolerance & 
Greater complexity, network latency & 
Massive data processing, scalable systems \\
\hline
\end{tabular}
\end{center}


Additional recommendations:
\begin{itemize}
    \item \textbf{\textcolor{accentColor}{Worker Threads:}} For CPU-intensive operations on a single machine (data processing, calculations).
    
    \item \textbf{\textcolor{accentColor}{Cluster Module:}} For I/O-bound applications needing to utilize all CPU cores (API servers).
    
    \item \textbf{\textcolor{accentColor}{Distributed Queues:}} When workloads exceed single-machine capacity or require fault tolerance.
    
    \item \textbf{\textcolor{accentColor}{Combined Approach:}} Use worker threads within each node of a distributed system for maximum parallelism.
\end{itemize}

\section{Conclusions}

Node.js's parallel and distributed computing capabilities effectively overcome its single-threaded limitations. The techniques presented enable developers to build computationally intensive applications that fully utilize modern multi-core architectures and distributed systems.

As data processing requirements continue to grow, these approaches become increasingly critical for maintaining performance and scalability in modern Node.js applications. The choice of the appropriate technique will depend on the type of workload and available resources: Worker Threads for CPU-intensive tasks on a single machine, Cluster for I/O-oriented applications, and distributed queues for systems requiring horizontal scalability.

\section{References}

\begin{itemize}
    \item Node.js. (2024). \textit{Worker Threads | Node.js v20.x Documentation}. \href{https://nodejs.org/docs/latest-v20.x/api/worker_threads.html}{Link}
    
    \item Node.js. (2024). \textit{Cluster | Node.js v20.x Documentation}. \href{https://nodejs.org/docs/latest-v20.x/api/cluster.html}{Link}
    
    \item Bull. (2023). \textit{Bull: Premium Queue Package for Node.js}. \href{https://github.com/OptimalBits/bull}{Link}
    
    \item Node.js Foundation. (2024). \textit{Event Loop and Performance}. \href{https://nodejs.org/en/docs/guides/event-loop-timers-and-nexttick/}{Link}
    
    \item Teixeira, P. (2023). \textit{Distributed Systems with Node.js}. O'Reilly Media.

    \item Translated, Edited and written in collaboration with AI.
\end{itemize}

\section{Explore My Other Posts}
\vspace{10pt}
\begin{tikzpicture}
  % Section background - increased height slightly to ensure all content fits
  \fill[color=terminalBg, rounded corners=5pt] (0,0) rectangle (\textwidth,-5);
  
  % Calculate width available for text (total width minus QR code width minus margins)
  \pgfmathsetmacro{\availableWidth}{\textwidth-4cm}
  
  % Left column content
  \begin{scope}
    % Section title
    \node[text=accentColor, font=\Large\bfseries, anchor=north west] at (0.5cm,-0.5cm) {Enjoyed This Content?};
    
    % Subtitle
    \node[text=primaryColor, font=\normalsize, anchor=north west] at (0.5cm,-1.2cm) {Don't miss my previous post about:};
    
    % Node.js Streams Part 1 title - with highlighted style and limited width
    \node[text=secondaryColor, font=\large\bfseries, anchor=north west, text width=\availableWidth] at (0.5cm,-1.9cm) {\color{accentColor}Node.js Streams: \color{primaryColor}Part 1 - Introduction \& Memory Efficiency};
    
    % Brief description of the post - with limited width and adjusted position
    \node[text=jsSecondary, font=\normalsize, text width=\availableWidth, anchor=north west] at (0.5cm,-3.1cm) {
      Learn the fundamentals of Node.js Streams and discover how they can dramatically reduce memory usage when processing large files.
    };
  \end{scope}
  
  % QR Code on the right - vertically centered
  \node[anchor=center] at ({\textwidth-1.5cm}, {-2.5cm}) {
    \qrcode[height=2.2cm]{https://www.linkedin.com/feed/update/urn:li:activity:7309333697692946432/}
  };
  
\end{tikzpicture}
\vspace{5pt}


\clearpage
\thispagestyle{empty}
\finalpagecontents

\end{document}